% Synthetic test file for the ECEB style linter (scripts/lint_euclid_style.py).
%
% Each violation line is marked with:    % EXPECT: <rule_id>
% Each clean counterpart is marked with: % CLEAN: <rule_id>
% Document-level expectations use:       % EXPECT_DOC: <rule_id>
% Edge cases use:                        % EDGE: <label>
%
% The linter should flag every EXPECT line and skip every CLEAN line.

% ====================================================================
% PREAMBLE — must be entirely skipped by the linter
% ====================================================================
\documentclass{article}
\usepackage{amsmath}
\usepackage{hyperref}
\newcommand{\Euclid}{\textit{Euclid}}
\newcommand{\IE}{$I_{\mathrm{E}}$}
\newcommand{\YE}{$Y_{\mathrm{E}}$}
\newcommand{\JE}{$J_{\mathrm{E}}$}
\newcommand{\HE}{$H_{\mathrm{E}}$}
% Preamble must NOT trigger rules even with violations present:
% dataset color analyze gray percent favor catalog comprised of
\definecolor{gray}{HTML}{808080}
\definecolor{lightgray}{gray}{0.9}

\begin{document}

% ====================================================================
% NAMING & TERMINOLOGY (N01--N10)
% ====================================================================

The Euclid mission observed the deep fields. % EXPECT: N01
The \Euclid\ mission observed the deep fields. % CLEAN: N01

The \textit{Euclid} Consortium approved the results. % EXPECT: N02
The Euclid Consortium approved the results. % CLEAN: N02

The \textsc{VIS} instrument was calibrated. % EXPECT: N03
The VIS instrument was calibrated. % CLEAN: N03

Observations were taken in the IE band. % EXPECT: N04
Observations were taken in the \IE\ band. % CLEAN: N04

We used the full dataset for training. % EXPECT: N05
We used the full data set for training. % CLEAN: N05

The sample is comprised of 500 galaxies. % EXPECT: N06
The sample is composed of 500 galaxies. % CLEAN: N06

The catalogue is publically available. % EXPECT: N07
The catalogue is publicly available. % CLEAN: N07

We achieved a high S/N ratio in all bands. % EXPECT: N08
We achieved a high S/N in all bands. % CLEAN: N08

The modelisation of the PSF is critical. % EXPECT: N09
The modelling of the PSF is critical. % CLEAN: N09

The errors associated to this method are small. % EXPECT: N10
The errors associated with this method are small. % CLEAN: N10

This will allow to detect faint sources in the field. % EXPECT: N11
This will allow one to detect faint sources in the field. % CLEAN: N11

The analysis uses point like sources as calibrators. % EXPECT: N12
The analysis uses point-like sources as calibrators. % CLEAN: N12

The mean redshift is $<z>\sim 1.6$ in the sample. % EXPECT: N13
The mean redshift is $\langle z \rangle \sim 1.6$ in the sample. % CLEAN: N13
The mean redshift is $\ave{z} \sim 1.6$ in the sample. % CLEAN: N13

The shear sample satisfies $N_s >> N_z$ in this regime. % EXPECT: N14
The shear sample satisfies $N_s \gg N_z$ in this regime. % CLEAN: N14
The faint sample has $L << L_*$ at this magnitude. % EXPECT: N14
The faint sample has $L \ll L_*$ at this magnitude. % CLEAN: N14
The constraint is $x < y$ and $a > b$ in all bins. % EDGE: N14_arrow
The threshold is $a <= b$ and $c >= d$ throughout. % EDGE: N14_eq

We compute the two point % EXPECT: N15
correlation function for the survey.
We compute the two-point correlation function for the survey. % CLEAN: N15

We obtain this via a Fourier Transform of the data. % EXPECT: N16
We obtain this via a Fourier transform of the data. % CLEAN: N16
The signal also requires a Fourier % EXPECT: N16
Transform across the data cube.

% ====================================================================
% BRITISH ENGLISH (E01--E08)
% ====================================================================

We analyze the photometric redshifts. % EXPECT: E01
We analyse the photometric redshifts. % CLEAN: E01

This accounts for 30 percent of the sample. % EXPECT: E02
This accounts for 30 per cent of the sample. % CLEAN: E02

The background appears gray in the image. % EXPECT: E03
The background appears grey in the image. % CLEAN: E03

See the acknowledgment section below. % EXPECT: E04
See the acknowledgement section below. % CLEAN: E04

Accurate modeling requires good data. % EXPECT: E05
Accurate modelling requires good data. % CLEAN: E05

Each source is labeled with a unique ID. % EXPECT: E06
Each source is labelled with a unique ID. % CLEAN: E06

We cross-matched the catalog with SDSS. % EXPECT: E07
We cross-matched the catalogue with SDSS. % CLEAN: E07

Results favor a flat universe. % EXPECT: E08
Results favour a flat universe. % CLEAN: E08

% ====================================================================
% UNITS & NUMBERS (U01, U02, U03, U05, U07, U08, U09, U10)
% ====================================================================

The velocity dispersion is 300 km/s in the core. % EXPECT: U01
The velocity dispersion is 300\,km\,s$^{-1}$ in the core. % CLEAN: U01

Separations were measured in arcsecs from the centre. % EXPECT: U02
Separations were measured in arcsec from the centre. % CLEAN: U02

The resolution is 5arcsec in all bands. % EXPECT: U03
The resolution is 5\,arcsec in all bands. % CLEAN: U03

The luminosity is 3e5 solar luminosities. % EXPECT: U05
The luminosity is $3 \times 10^{5}$ solar luminosities. % CLEAN: U05

The catalogue contains 10,000 sources. % EXPECT: U07
The catalogue contains 10\,000 sources. % CLEAN: U07

We expect to detect 100000 galaxy clusters in this survey. % EXPECT: U08
We expect to detect 100\,000 galaxy clusters in this survey. % CLEAN: U08
The dataset contains 1234567 individual measurements. % EXPECT: U08
The 2024 release supersedes the 2018 catalogue. % EDGE: U08_year
We targeted SDSS J012345 in the deep field. % EDGE: U08_id

The survey covers 1\,900 square degrees of sky. % EXPECT: U10
The reference sample contains $4\,352$ stars. % EXPECT: U10
In DR1 there are more than 3 000 exposures. % EXPECT: U10
The catalogue lists 1{,}900 bright stars. % EXPECT: U10
The sample has 1,900 stars in total. % CLEAN: U07
The survey covers 1900 square degrees of sky. % CLEAN: U10
There are 42\,092 calibrated frames in DR1. % CLEAN: U10
We processed \num{4217} individual frames. % CLEAN: U10
In Sect. 4 300 sources were catalogued. % CLEAN: U10

The galaxy density is $5\, 10^7$ per square degree. % EXPECT: U09
The galaxy density is $5\cdot 10^7$ per square degree. % EXPECT: U09
The galaxy density is $5\times 10^7$ per square degree. % CLEAN: U09

% ====================================================================
% TYPESETTING (T01, T02, T04, T05, T06, T08, T09, T10, T11, T12)
% ====================================================================

The so-called "standard" candles were observed. % EXPECT: T01
The so-called ``standard'' candles were observed. % CLEAN: T01

We found (20-25) candidate sources per field. % EXPECT: T02
We found (20--25) candidate sources per field. % CLEAN: T02

The fit gives $log(L) = 42.3$ for the sample. % EXPECT: T04
The fit gives $\log(L) = 42.3$ for the sample. % CLEAN: T04

This paragraph ends with a forced line break.\\ % EXPECT: T05
This paragraph ends normally. % CLEAN: T05

See the documentation at https://euclid.esa.int for details. % EXPECT: T06
See the documentation at \url{https://euclid.esa.int} for details. % CLEAN: T06

Fig. 1 shows the photometric residuals. % EXPECT: T08
Figure 1 shows the photometric residuals. % CLEAN: T08

\includegraphics[width=0.8\columnwidth, height=4cm]{fig.pdf} % EXPECT: T09
\includegraphics[width=0.8\columnwidth]{fig.pdf} % CLEAN: T09

The sources (AGN and stars)(see Merloni et al. 2018) are detected. % EXPECT: T10
The sources (AGN and stars; see Merloni et al. 2018) are detected. % CLEAN: T10

\acknowledgement{We thank the referee.} % EXPECT: T11
\begin{acknowledgements} % CLEAN: T11

This can be expressed as:
\begin{equation} % EXPECT: T12
x = y + z
\end{equation}

This can be expressed as
\begin{equation} % CLEAN: T12
x = y + z
\end{equation}

This was the first ''Stage-IV'' dark energy experiment. % EXPECT: T13
This was the first ``Stage-IV'' dark energy experiment. % CLEAN: T13
The so-called ``standard'' candles were observed. % EDGE: T13_closer

The detector is sensitive at $E=1.5keV$ in this band. % EXPECT: T14
The detector is sensitive at $E=1.5\,\mathrm{keV}$ in this band. % CLEAN: T14
The line is at $\nu=1.4GHz$ for HI emission. % EXPECT: T14
The size is about $1Mpc$ across the cluster. % EXPECT: T14

% ====================================================================
% REFERENCES & CITATIONS (R02, R03, R04, R05)
% ====================================================================

Euclid Collaboration et al. (2024) reported the calibration. % EXPECT: R02
Euclid Collaboration: Scaramella et al. (2024) reported the calibration. % CLEAN: R02

% This old paragraph was part of an earlier draft and should have been deleted before the final submission % EXPECT: R03
% Short note % CLEAN: R03

Smith, J. 2023, arXiv e-prints, arXiv:2301.12345 % EXPECT: R05
Smith, J. 2023, arXiv:2301.12345 % CLEAN: R05

% R04 is a document-level check: absence of the AckEC macro means it fires % EXPECT_DOC: R04
% R06 (release=dr1, default): absence of AckDRone/DR1cite means it fires % EXPECT_DOC: R06

% ====================================================================
% STYLE GUIDE SPECIFIC (S01--S05)
% ====================================================================

The DEC of the target is measured precisely. % EXPECT: S01
The Dec of the target is measured precisely. % CLEAN: S01

The distribution is nonGaussian in the tails. % EXPECT: S02
The distribution is non-Gaussian in the tails. % CLEAN: S02

Deep imaging in the J band reached 26 mag. % EXPECT: S03
Deep imaging in the $\textit{J}$ band reached 26 mag. % CLEAN: S03

The data is consistent with earlier measurements. % EXPECT: S04
The data are consistent with earlier measurements. % CLEAN: S04

The expansion of the universe is accelerating. % EXPECT: S05
The expansion of the Universe is accelerating. % CLEAN: S05

% ====================================================================
% EDGE CASES — must NOT produce false positives
% ====================================================================

% Edge: \definecolor{gray} in body must NOT trigger E03/E01
\definecolor{mygray}{gray}{0.5} % EDGE: E03-definecolor

% Edge: \newcommand with Euclid must NOT trigger N01
\newcommand{\MyEuclid}{\textit{Euclid}} % EDGE: N01-newcommand

% Edge: \url{} wrapping must NOT trigger T06
Visit \url{http://example.com/data} for details. % EDGE: T06-url

% Edge: math content must NOT trigger U01
The velocity is $v = 300\,\mathrm{km\,s^{-1}}$ in the cluster. % EDGE: U01-math

% Edge: \texttt{} must NOT trigger E01
Use the \texttt{color} keyword in the configuration file. % EDGE: E01-texttt

% Edge: verbatim environment must NOT trigger N05
\begin{verbatim}
dataset = load("input.csv")
color = "gray"
\end{verbatim}

% Edge: display math (equation env) must NOT trigger T04
\begin{equation}
\log_{10} x = y
\end{equation}

% Edge: nested environments (aligned inside equation inside figure)
\begin{figure}
\begin{equation}
\begin{aligned}
  a &= b + c \\
  d &= e \times f
\end{aligned}
\end{equation}
\caption{Test figure.}
\end{figure}

% Edge: \ac{} macros must NOT trigger N01
The \ac{Euclid} mission is described in detail. % EDGE: N01-ac

% Edge: \href{url}{text} must NOT trigger T06
See \href{https://euclid.esa.int}{the website} for more. % EDGE: T06-href

% Edge: inline math with \, before unit must NOT trigger U03
The distance is $5\,\mathrm{arcsec}$ from the centre. % EDGE: U03-math

% Edge: number range inside \ref must NOT trigger T02
See Table~\ref{tab2-5} for a summary. % EDGE: T02-ref

% Edge: \textit{Euclid} (correct usage) must NOT trigger N01
The \textit{Euclid} satellite was launched in 2023. % EDGE: N01-textit

% Edge: R.A. test for S01
The R.A. of the cluster is measured. % EXPECT: S01

% Edge: ORCID line must NOT trigger T02
\and J.~Smith\orcid{0000-0002-5258-8761}\inst{\ref{aff3}} % EDGE: T02-orcid

% Edge: affiliation address must NOT trigger T02
Kavli Institute, Chiba 277-8583, Japan\label{aff103} % EDGE: T02-address

% Edge: {\tt Euclid Emulator} must NOT trigger N01
The {\tt Euclid Emulator 2} reproduces nonlinear clustering. % EDGE: N01-tt

% Edge: "Center for" (institution name) must NOT trigger E01
Center for Computational Astrophysics, New York. % EDGE: E01-institution

% Edge: \begin{center} must NOT trigger E01
\begin{center}
This text is centred.
\end{center}

% Edge: postal code must NOT trigger U05
NRC Herzberg, Victoria, BC V9E 2E7, Canada\label{aff224} % EDGE: U05-postal

% Edge: "allow" with an object must NOT trigger N11
This will allow us to detect faint sources. % EDGE: N11-object

% Edge: hyphenated compound must NOT trigger N12
The analysis uses point-like sources as calibrators. % EDGE: N12-hyphenated

% Edge: \includegraphics with only width must NOT trigger T09
\includegraphics[width=0.8\columnwidth]{fig.pdf} % EDGE: T09-width-only

% Edge: separated parentheses with text between must NOT trigger T10
The sources (AGN and stars) and (see Merloni et al. 2018) are detected. % EDGE: T10-separated

% Edge: equation without colon on prev line must NOT trigger T12
This can be expressed as
\begin{equation} % EDGE: T12-no-colon
x = y + z
\end{equation}

% Edge: "a different universe" must NOT trigger S05
In a different universe the laws of physics might differ. % EDGE: S05-different

% Edge: "the galaxy" (generic, not Milky Way) must NOT trigger S05
The galaxy has a stellar mass of $10^{11}$ solar masses. % EDGE: S05-generic-galaxy

% Edge: affiliation \\ must NOT trigger T05
$^{1}$ INAF-Osservatorio, Via Frascati 33, I-00078 Roma, Italy\\ % EDGE: T05-affiliation

% Edge: postal codes inside a multi-line \institute{...} block must NOT
% trigger U08 (thousands separator), even on continuation lines that carry
% no line-local affiliation marker.
\institute{Institut d'Astrophysique de Paris, 98bis bd Arago, 75014 Paris, France\label{aff301} % EDGE: U08-zip-institute
\and
Max-Planck-Institut f\"ur extraterrestrische Physik, Giessenbachstra{\ss}e 1, % EDGE: U08-zip-wrap-head
85748 Garching, Germany\label{aff302} % EDGE: U08-zip-wrapped
\and
Univ. Grenoble Alpes, CNRS, IPAG, 38000 Grenoble, France % EDGE: U08-zip-block-only
\and
Dipartimento di Fisica e Astronomia "G. Galilei", Universit\`a di Padova, 35131 Padova, Italy\label{aff303} % EDGE: T01-affiliation-quotes
}
The survey contains 20000 sources in this field. % EXPECT: U08

% Edge: affiliation-reference list in \inst{...} must NOT trigger U07
\and C.~Author\inst{12,345} % EDGE: U07-inst-refs

% Edge: underscore-joined identifier must NOT trigger N01
The transient was internally called Euclid\_SNT\_2023B in the pipeline. % EDGE: N01-identifier

% Edge: LaTeX dimension arguments must NOT trigger T14
\vspace{-2.5cm} % EDGE: T14-vspace
\rule{0pt}{3cm} % EDGE: T14-rule

% Edge: \\ separating entries inside \tablefoot{} must NOT trigger T05
\tablefoot{These fiducial values are used for the synthetic data.
\\ % EDGE: T05-tablefoot
$^\ddag$ A constant nuisance parameter is assumed.}

% Edge: "data" as infinitive object must NOT trigger S04
The software that is employed to process the data is under strict control. % EDGE: S04-infinitive-object

% Edge: reference abbreviations before \ref are not sentence ends (T15)
Details are given in Sect. \ref{sec:model} and in Fig. \ref{fig:layout}. % EDGE: T15-abbrev-midsentence
The dither pattern is shown in Fig. % EDGE: T15-abbrev-wrap-head
\ref{fig:dither} for the reference sequence. % EDGE: T15-abbrev-linestart

% Edge: title-case compound names keep roman Euclid (N01)
The Euclid Reference Survey Definition was produced by the Euclid Science Team. % EDGE: N01-titlecase-chain
WISHES stands for Wide Imaging with Subaru HSC of the Euclid Sky. % EDGE: N01-sky
\bibliography{refs,Euclid} % EDGE: N01-bibliography

% Edge: proper-name Catalog is not anglicised (E07)
Stars come from the 2MASS Point Source Catalog for calibration. % EDGE: E07-proper-name

% Edge: communications bands are not photometric wavebands (S03)
The spacecraft uses K-band telemetry for the downlink. % EDGE: S03-telemetry

% Edge: label/ref bodies are not prose (U03)
\label{fig:deepn-20deg-area-time} % EDGE: U03-label-body
The survey area is shown in Fig.~\ref{fig:deepn-20deg-area-time} here. % EDGE: U03-ref-body

% Edge: table constructs (T14 column spec / \vskip, U10 alignment, T02 \cline)
\begin{tabular}{l p{6cm}} % EDGE: T14-colspec
\noalign{\vskip 4mm} % EDGE: T14-vskip-noalign
Pass & \phantom{0}2\,656 \\ % EDGE: U10-table-alignment
\cline{1-2} % EDGE: T02-cline
Total & 12\,673 \\
\end{tabular}

% Edge: standalone noun "power law" must NOT trigger N12
The relation is described by a power law with slope $\eta$. % EDGE: N12-noun

% Edge: Euclid at the tail of a compound name must NOT trigger N01
We use the Cosmology Likelihood for Observables in Euclid pipeline. % EDGE: N01-tail
The Likelihood Tool for Euclid analyses are presented. % EDGE: N01-tail-tool

% Edge: abbreviation on wrapped line (no sentence break) must NOT trigger T08
The residuals are shown in
Fig. 1 as red points. % EDGE: T08-wrapped

% Genuine sentence start across lines: prev line ends with "."
The analysis is complete.
Fig. 1 shows the results. % EXPECT: T08

% Edge: line after "et al." should NOT trigger T08 (ends in "." but not a sentence)
As described by Smith et al.
Fig. 2 illustrates the trend. % EDGE: T08-etal

% Edge: catalogue/star name with HE+digits must NOT trigger N04
The spectrum of HE 0107-5240 is well studied. % EDGE: N04-star

% Edge: "dialog box" (software) must NOT trigger E01
The dialog box prompts for a filename. % EDGE: E01-dialog

% T04 in display math should be detected
\begin{equation}
y = log(x) + sin(x) % EXPECT: T04
\end{equation}

% T04 in \[...\] should be detected
\[ z = cos(\theta) \] % EXPECT: T04

% R03: short non-prose comment must NOT trigger
% See note.
% R03: long prose comment SHOULD trigger
% We have removed this paragraph because the referee did not like it.

% Edge: "Euclid" inside \includegraphics{} filename must NOT trigger N01
\includegraphics[width=0.5\textwidth]{path/Euclid-image.png} % EDGE: N01-graphics-path
\includegraphics{Figures/Euclid.pdf} % EDGE: N01-graphics-bare

% Edge: mixed-case catalogue prefix (Abell, Hickson) must NOT trigger T02
The cluster Abell 1689-23 is a strong lens. % EDGE: T02-abell
Hickson 92-3 is a compact group. % EDGE: T02-hickson

% Edge: \num{} from siunitx must NOT trigger U08
We detect \num{15000} galaxies in the deep field. % EDGE: U08-num
The catalogue contains \num{1234567} sources. % EDGE: U08-num-large

% ---- N17: Gaia (mission) should be italicised (mirrors N01 for Euclid) ----
We cross-match VIS detections with Gaia DR3 sources. % EXPECT: N17
The \Gaia\ mission delivered its third data release. % CLEAN: N17
The Gaia Collaboration published the catalogue. % CLEAN: N17
\includegraphics{figures/Gaia_map.pdf} % CLEAN: N17

% ---- T15: lowercase \cref/\ref at sentence start -> use \Cref -------------
The trend is clear. \cref{fig:demo} shows the residuals. % EXPECT: T15
The measurement converges after several iterations.
\cref{fig:demo3} summarises the outcome. % EXPECT: T15
As shown in \cref{fig:demo4}, the residuals are flat. % CLEAN: T15
\Cref{fig:demo5} shows the map. % CLEAN: T15
Results agree with Smith et al. \cref{fig:demo6} shows the fit. % CLEAN: T15

% ---- N01/N02: Euclid data-product names stay roman ------------------------
The Euclid \ac{DR1} products are now public. % CLEAN: N01
Delivered as part of the Euclid \acl{DR1} in 2026. % CLEAN: N01
The VIS Euclid DR1 catalogues are matched to references. % CLEAN: N01
We study galaxies in the Euclid Ultra-Deep Field. % CLEAN: N01
The \Euclid DR1 images are now public. % EXPECT: N02
The \Euclid\ Q1 data were released earlier. % EXPECT: N02
The \textit{Euclid} \ac{DR1} frames are calibrated. % EXPECT: N02
The \Euclid \ac{VIS} camera observes one field at a time. % CLEAN: N02
Positions from \textit{Gaia} DR3 are used throughout. % CLEAN: N17

% ---- T16: panel descriptors in captions -> \emph{Left}: -------------------
\begin{figure}
\caption{Left: the model prediction. Right: the residuals.} % EXPECT: T16
\end{figure}
\begin{figure}
\caption{\emph{Left:} the model prediction.} % EXPECT: T16
\end{figure}
\begin{figure}
\caption{\emph{Left}: the model. \textit{Right panel}: the residuals.} % CLEAN: T16
\end{figure}
\begin{figure}
\caption{Arrows on the right panel are enlarged four times.} % CLEAN: T16
\end{figure}
\begin{figure}
\caption{Focal-plane residuals over all observations.
  Lower right: residuals after subtraction.} % EXPECT: T16
\end{figure}
Left: this is body prose, not a caption. % EDGE: T16-prose-not-caption
\begin{figure}
\caption{The upper limit: derived at 95 per cent confidence.} % EDGE: T16-upper-limit
\end{figure}

% ---- U08/T02: identifiers are not quantities or ranges --------------------
The standard has SourceID 1440758225532172032 in the archive. % EDGE: U08-gaia-source-id
Object 12345678901 lies in the deep field. % EDGE: U08-long-integer
The star has SourceID 98765 in the list. % EDGE: U08-source-id
The source has IDENT 271158 in the catalogue. % EDGE: U08-ident
Ground testing for \ac{CCD} 6-2 shows nonlinearity. % EDGE: T02-ccd-space

% ---- T17: blank line after equation, sentence continues lowercase ---------
\begin{equation}
  a = b + c
\end{equation}

where $b$ is held fixed. % EXPECT: T17
\begin{equation}
  d = e
\end{equation}
where $e$ is free. % CLEAN: T17
\begin{equation}
  f = g
\end{equation}
%
where $g$ is bound. % CLEAN: T17
\begin{equation}
  h = k
\end{equation}

The next paragraph opens correctly. % CLEAN: T17

\end{document}
